\chapter{Open Problems under the Three-Axis Discipline}
\label{ch:open-frontiers}

\providecommand{\Zder}{Z^{\mathrm{der}}_{\mathrm{ch}}}
\providecommand{\GRTone}{\mathrm{GRT}_1(\mathbb{Q})}

The discipline of Appendix~\ref{app:three-axis-scope-discipline} separates
proved comparison maps from conditional ones. The open problems below are
the remaining comparison statements after scope-omission collapses are removed.
Each carries explicit
$(\text{level}, \text{chart}, \text{ambient})$ coordinates, its hypotheses,
and the reconstruction theorem to be proved.

\section{The chain fusion conjecture in general $d$}
\label{sec:open-chain-fusion}

\paragraph{Coordinates.} Level 2 (Stage-2 chiral algebra equals open-side
boundary algebra). Chart $(\Sigma_{d-1}, C, b, \mathrm{adm})$ with
$(C, D_C, \tau_C)$ the tangential log curve on the open side. Ambient:
chain-level on the constructed loci and $(\infty, 1)$-categorical for the
general statement.

\paragraph{Hypotheses.}
\begin{itemize}[nosep]
\item $X$ a $\mathrm{CY}_d$-target, $\cC_X$ the cyclic $\Ainf$-CY
  category at level 0.
\item $(\Sigma_{d-1}, C)$ a $(d{-}1)$-cycle in $X$ specialised to a
  reference curve $C$.
\item $(C, D_C, \tau_C)$ a tangential log curve on $C$ with $D_C$
  encoding the special points of the CY datum (orbifold loci, fibration
  punctures, conifold singularities); the definition of $(X, D, \tau)$
  is the one of Vol~I,
  \cite{VolI}.
\item An open factorisation dg-category on $(C, D_C, \tau_C)$ in the
  Vol~I/II sense, with closed-colour input
  $(\cC_X^{\mathrm{op}}, \Theta_{\cC}, \mathrm{Tr}_{\cC})$.
\item A canonical boundary vacuum $b(X, \Sigma, C)$ in this open
  factorisation category, induced by the CY data.
\end{itemize}

\paragraph{Conjecture.} For every CY$_d$-category
$\cC_X$ and every chart $(\Sigma_{d-1}, C)$, the Stage-2 chiral algebra
\[
A_X \;=\; \Phi^{(\Sigma_{d-1}, C)}_d(\cC_X)
\;=\; \SpCh_{\Sigma_{d-1}, C} \!\circ \PhiFA_d(\cC_X)
\quad\text{on } C
\]
agrees with the boundary algebra
\[
A_{b(X, \Sigma, C)} \;=\; \mathrm{End}^{\mathrm{op-fact}}_{\cC_X}\bigl(b(X, \Sigma, C)\bigr)
\]
of the open factorisation category at the canonical boundary vacuum.

\paragraph{Constructed loci.} $\mathbb{C}^3$ via Costello--Li 6d hCS;
local $\mathbb{P}^2$ via the Beilinson 9-arrow McKay quiver;
resolved conifold via the Szendr\H{o}i 2-vertex Czech atlas;
$K3 \times E$ via the K\"unneth reduction
$D^b\mathrm{Coh}(K3 \times E) \simeq D^b\mathrm{Coh}(K3) \otimes D^b\mathrm{Coh}(E)$
plus the Costello--Gaiotto--Yagi 5d hCS Yangian VOA on the elliptic
factor. All four cases construct $b(X, \Sigma, C)$ explicitly and
verify the Stage-2-equals-boundary identification chain-level.

\paragraph{Construction.}
\begin{enumerate}[nosep]
\item \emph{Boundary vacuum existence.} For each equivariance stratum
  (toric / reduced + Aut / orbifold inertia / lattice-polarised), state
  the canonical-vacuum construction. The toric case admits the
  perverse coherent-sheaf vacuum sourced by the GIT chamber framing;
  the reduced + Aut case admits the K\"unneth-decomposed vacuum on the
  fibre product; the orbifold inertia case admits the
  $I(X/G)$-equivariant vacuum; the lattice-polarised case admits the
  period-domain vacuum. Stratum-by-stratum existence is a finite
  enumeration of constructions.
\item \emph{Stage-2 identification.} Given $b(X, \Sigma, C)$, identify
  the open-factorisation algebra $\mathrm{End}^{\mathrm{op-fact}}(b)$
  with the closed Stage-2 image $A_X$ via the swiss-cheese pair
  comparison of Vol~II (boundary $A$, universal closed sector
	  $\Zder(A)$; physical-bulk reading conditional on the OCA arrow
  of Definition~\ref{def:oca-quasi-iso}).
\item \emph{Compactness obstruction.} Compact non-toric CY$_3$ is the
  hardest case: the open factorisation category requires a tangential
  log curve $(C, D_C, \tau_C)$ derived from the period datum, and the
  CY-A$_3$ chain-level closure (\S\ref{sec:open-cya3-chain})
  feeds the construction. Resolution of CY-A$_3$ chain-level closes
  this obstruction.
\item \emph{Universality.} The collection
  $\{\Phi^{(\Sigma_{d-1}, C)}_d\}_{d, (\Sigma, C)}$ is a two-parameter
  correspondence family (parameters $d$ and chart datum). Stage-1
  functoriality plus Stage-2 chart-specialisation gives functoriality
  on objects; functoriality on morphisms reduces to per-$d$
  $R$-matrix-gauge transport.
\end{enumerate}

\paragraph{Obstruction.}
\begin{itemize}[nosep]
\item \emph{Counterexample check.} A non-formal compact CY$_3$ with no
  global Lagrangian boundary brane could obstruct the canonical-vacuum
  construction. At compact non-toric CY$_3$ the obstruction is the
  absence of a canonical boundary vacuum outside the constructed loci
  of \S\ref{sec:open-Gx-compact}.
\item \emph{Equivalence check.} ``Equals $A_{b(X, \Sigma, C)}$'' means
  quasi-equivalence of $E_1$-chiral algebras on $C$, with the
  equivalence implemented by an explicit chain map. The four verified
  loci provide the chain map; the general case is conjectural.
\end{itemize}

\section{$G(X) = D(Y^+(X))$ for compact non-toric CY$_3$}
\label{sec:open-Gx-compact}

\paragraph{Coordinates.} Level 3 (positive half $\to$ Drinfeld double).
Chart: compact non-toric CY$_3$ with reduced $\mathbb{C}^\times +$ Aut
or lattice-polarised stratum. Ambient: chain-level with
inverse-limit / pro-completion.

\paragraph{Hypotheses.} The conditions:
\begin{itemize}[nosep]
\item Compact critical CoHA construction on $\cM^+_{\mathrm{eff}}(X)$
  with $\phi_W$ vanishing-cycle sheaf.
\item Hall--Drinfeld doubling: extension from positive half $Y^+(X)$ to
  the full Drinfeld double $G(X) = D(Y^+(X))$.
\item Radical descent: identification of the radical of the Hall pairing
  on the compact critical CoHA.
\item PBW with no extra relations: triangular decomposition of $G(X)$
  with positive, negative, and Cartan halves of the expected dimension.
\item Parity (Hodge-parity $\mathbb{Z}/2$-grading): consistent
  $\mathbb{Z}/2$-super-extension on the Mukai self-mirror branch.
\item Green-adjoint coproduct: extension of the bialgebra coproduct
  from toric / orbifold loci to compact non-toric.
\item Primitive-centre reduction: identification of the centre
  $\mathfrak{z}(G(X))$ with the level-3 chiral centre $\Zder(A_X)$.
\item Associator cohomology: Drinfeld associator class in $H^3$ of the
  compact non-toric Hall datum.
\item Completion / inverse-limit consistency: the inverse system
  $\{Y^+_N(X)\}$ over filtration $N$ admits a Mittag--Leffler condition
  and converges to a compact $Y^+(X)$.
\item Heegner comparison: Heegner-class restriction of the compact
  $Y^+(X)$ matches the period-domain Yangian datum.
\end{itemize}

\paragraph{Conjecture.} For compact non-toric
$\mathrm{CY}_3$ $X$ in the reduced $\mathbb{C}^\times +$ Aut or
lattice-polarised stratum, the positive half
$Y^+(X) = H^\bullet_{\mathrm{eq}}(\cM^+_{\mathrm{eff}}(X), \phi_W)$
exists as an associative $E_1$-algebra in the chain-level ambient, the
Drinfeld double $G(X) = D(Y^+(X))$ is a Hopf algebra, and the
identification
\[
\Zder(A_X) \;\simeq\; \mathfrak{z}\bigl(G(X)\bigr)
\]
holds chart-by-chart at the chain-level.

\paragraph{Constructed loci.} Toric ($\mathbb{C}^3$, local $\mathbb{P}^2$,
resolved conifold) via Schiffmann--Vasserot, Davison, Padurariu, Maulik--Okounkov.
Orbifold inertia $\mathbb{C}^3 / \Gamma$ for $\Gamma \subset \mathrm{SU}(3)$
finite via Bridgeland--King--Reid. K3 fibrations over rational base
partially via the elliptic-genus extension. Compact non-toric (quintic,
generic Schoen) is unconstructed.

\paragraph{Construction.}
\begin{enumerate}[nosep]
\item Construct compact critical CoHA on a Verlinde-type family
  (the quintic Fermat with $\mathbb{Z}_5^5$-orbifold action) and verify
	  the conditions one-by-one. The $\mathbb{Z}_5^5$ symmetry reduces the
  computation to $5^5$-equivariant cohomology of the quotient stack.
\item Heegner-comparison condition as the cleanest single criterion:
  identification of the Heegner-class divisors on the period-domain
  with the radical of the Hall pairing on the chosen family.
\item Stage from $K3 \times E$ (verified) through the elliptic
  fibration $X \to \mathbb{P}^1$ structure on the Schoen Calabi--Yau
  (compact non-toric, $\chi = 0$, fibred): the K\"unneth-multiplicative
  reduction at the K3 fibre extends through the elliptic fibration with
  the period-domain Yangian datum at the fibre.
\end{enumerate}

\paragraph{Obstruction.}
\begin{itemize}[nosep]
\item \emph{Missing hypothesis check.} The compactness condition
  (Mittag--Leffler condition) is the deepest. A compact non-toric
  inverse system that fails Mittag--Leffler would refute the existence;
  no counterexample is currently known.
\item \emph{Numerical-constant check.} In the computed degrees for
  \(K3\times E\), the chart-local generators see only a small quotient
  of the global Drinfeld centre.  The compact CoHA construction must
  therefore include the Maulik--Okounkov stable-envelope gluing class;
  chart-local Hall data alone cannot determine the centre.
\end{itemize}

\section{$W_\infty[\lambda] \Rightarrow E_\infty$ outside the
admissible window}
\label{sec:open-Winfty-Einfty}

\paragraph{Coordinates.} Level 3 (Fock evaluation endpoint at the
$\lambda$-chart). Chart: chosen Fock evaluation $\lambda$. Ambient:
admissibility window (a four-condition assembly).

\paragraph{Hypotheses.} The four-condition admissible window:
\begin{itemize}[nosep]
\item Pro\v{c}\'azka triangular truncation: the truncation
  $W_\infty[\lambda] \to W_N[\lambda]$ exists as a triangular
  decomposition.
\item Creutzig--Kanade--Linshaw parafermion construction: the
  parafermion coset
  $W_\infty[\lambda] / W_\infty^{\mathrm{ab}}[\lambda]$ is identified
  with a coset VOA.
\item Pope--Romans--Shen / Bakas presentation: the central charge
  $c(\lambda)$ admits the Pope--Romans--Shen quadratic dependence and
  the Bakas $\mathfrak{w}$-algebra structure.
\item Yamada weight window: the conformal weight spectrum lies in the
  Yamada window $\Delta \in [\Delta_{\min}, \Delta_{\max}]$ for
  $\Delta_{\min}, \Delta_{\max}$ depending on $\lambda$.
\end{itemize}

\paragraph{Conjecture.} The implication
$W_\infty[\lambda] \Rightarrow E_\infty$ holds inside the four-condition
admissible window. Outside the window, the implication is open: identify
the obstruction (which condition fails) and the next admissible window.

\paragraph{Verified subwindow.} Spin $\leq 8$ with $\lambda$ in the
Pro\v{c}\'azka domain (Bakas--Pope--Romans--Shen quadratic) is verified.
Full-spin and exotic $\lambda$ are open.

\paragraph{Construction.}
\begin{enumerate}[nosep]
\item Replace the Pope--Romans--Shen quadratic with Yamada's full
  weight spectrum analysis to extend the window to include exotic
  $\lambda$ with non-quadratic central-charge dependence.
\item Use the Costello--Gaiotto holomorphic factor analysis to extend
  the window past spin $8$, treating each spin-$s$ generator as an
  $E_\infty$-deformation of the spin-$(s-1)$ generator.
\item Identify the boundary of the admissible window with a sharp
  $W_\infty$-cohomological obstruction (proving outside the window the
  $E_\infty$-structure obstructs at named cohomological depth).
\end{enumerate}

\paragraph{Obstruction.}
\begin{itemize}[nosep]
\item \emph{Convergence check.} The implication is asserted in the
  weight-completed ambient; the chain-level statement in ordinary
  complexes is open and may genuinely fail.
\item \emph{Theory-import check.} The $E_\infty$-structure is the
  Costello $E_\infty$ on factorisation algebras; equivalences across
  different $E_\infty$-conventions (operadic versus simplicial) must be
  named.
\end{itemize}

\section{Modularity compatibility under chain fusion}
\label{sec:open-modularity-fusion}

\paragraph{Coordinates.} Level 4 (scalar / modular form) versus level 3
(operator algebra) under the chain fusion at level 2. Chart: choice of
$\Phi_N$ on the CY-side; choice of clutching torus on the open side.
Ambient: HS-sewing on the open side; period-domain on the CY side.

\paragraph{Hypotheses.}
\begin{itemize}[nosep]
\item Chain fusion at level 2 (\S\ref{sec:open-chain-fusion}).
\item Open-side modular trace + clutching at level 4 (Vol~II:
  modular trace on $\overline{M}_{g, n}$ with clutching maps; the
  closed shadow at level 4 has modular consequence on
  $\Sp_{2g}(\mathbb{Z})$).
\item CY-side $\kappa_{\mathrm{BKM}}(\Phi_N) = c_N(0)/2$ universal
  identity at the chosen Siegel input denominator $\Phi_N$.
\end{itemize}

\paragraph{Conjecture.} Under chain fusion at
level 2, the open-side trace + clutching scalar at level 4 equals the
CY-side $\kappa_{\mathrm{BKM}}(\Phi_N)$ for the matching chart input.
The two scalars at level 4, computed by two different routes,
coincide.

\paragraph{Constructed loci.} $K3 \times E$ at $\Phi_N = \Delta_5$:
$\kappa_{\mathrm{BKM}}(\Delta_5) = 5$ on the CY side; the K3 elliptic
genus on the open side gives the matching trace. The
genus-1 partition function $\mathrm{tr}_{V^\natural} q^{L_0 - c/24}$ on
the Borcherds Monster matches at $d = 3$. The Fake-Monster
$\Phi_{12}$ trace matching is open.

\paragraph{Construction.}
\begin{enumerate}[nosep]
\item Identify the open-side modular trace as the closed-sector
  scalar trace on $\Zder(A)$ via the swiss-cheese pair (boundary
  $A$, universal closed sector $\Zder(A)$; physical-bulk reading
	  conditional on the OCA arrow of Definition~\ref{def:oca-quasi-iso}).
\item Apply the Borcherds-product identity
  $\Phi_N(Z)^{-1} = \prod_{(n, \ell, m) > 0} (1 - q^n y^\ell p^m)^{-c_N(nm - \ell^2/4)}$
  on the CY side to obtain the matching scalar.
\item The matching scalar identity reduces to a Petersson-pairing
  identity between the Igusa modular form on the CY side and the
  Borcherds-lifted Jacobi form on the open side.
\end{enumerate}

\paragraph{Obstruction.}
\begin{itemize}[nosep]
\item \emph{Sign / convention check.} The orientation gerbe on the
  open side and the sign of the Borcherds product on the CY side must
  be matched chart-by-chart. The CHL ladder
  $N \in \{1, 2, 3, 4, 6\}$ has signs determined by $\varphi(N) \mid 2$.
\item \emph{False functoriality check.} The open-side modular trace is
  a pairing on $\overline{M}_{g, n}$, not a single number; matching to
  the CY-side scalar requires choice of clutching torus.
\end{itemize}

\section{The universal Borcherds-weight identity as fusion consequence}
\label{sec:open-borcherds-as-fusion}

\paragraph{Coordinates.} Level 4 from level 3 via the trace pairing on
the chiral centre. Chart: Siegel input denominator $\Phi_N$. Ambient:
chain-level on the period domain; pro-completion on the chiral centre.

\paragraph{Hypotheses.}
\begin{itemize}[nosep]
\item Chain fusion at level 2 (\S\ref{sec:open-chain-fusion}).
\item Existence of the chiral centre $\Zder(A_X)$ at level 3
  (\S\ref{sec:open-Gx-compact} for compact non-toric).
\item Trace pairing $\mathrm{Tr}_{\Zder(A_X)} : \Zder(A_X) \to k[-d]$
  inherited from the CY trace on $\cC_X$.
\item Borcherds product identity at the chosen Siegel input $\Phi_N$.
\end{itemize}

\paragraph{Conjecture.} The universal
Borcherds-weight identity
\[
\kappa_{\mathrm{BKM}}(\Phi_N) \;=\; c_N(0)/2
\]
is expected to be the trace of an explicit operator at the level-3
chiral centre \(\Zder(A_X)\) at the chart \(\Phi_N\), but only after the
operator and the trace pairing have been constructed from the CY trace
at level \(0\) through the named comparison arrows.

\paragraph{Constructed scalar loci.} At the K3 denominator,
\[
 \kappa_{\mathrm{BKM}}(\Delta_5)=10/2=5.
\]
At the Fake-Monster denominator,
\[
 \kappa_{\mathrm{BKM}}(\Phi_{12})=24/2=12.
\]
These are level-$4$ Borcherds-weight identities.  Their interpretation
as traces on \(\Zder(A_{K3\times E})\), on the conditional \(d=5\)
Fake-Monster centre, or on the CHL-twisted centres requires the compact
Hall source, Hopf pairing, PBW/parity data, and centre trace comparison.

\paragraph{Construction.}
\begin{enumerate}[nosep]
\item Construct $\Zder(A_X)$ on the four equivariance strata
  (\S\ref{sec:open-Gx-compact}).
\item Compute the Cartan-generator trace at the level-3 chiral centre
  in the $\Phi_N$ chart, using the equivariant cohomology presentation
  of $Y^+(X)$.
\item Identify the trace with the Borcherds Fourier coefficient
  $c_N(0)/2$ via the Borcherds product identity at $\Phi_N$, applying
  Bruinier's Heegner Chern-class reciprocity for the lattice-polarised
  case.
\end{enumerate}

\paragraph{Obstruction.}
\begin{itemize}[nosep]
\item \emph{Independence check.} The universal Borcherds-weight identity
  may be a level-4-internal identity (Borcherds 1995 \emph{Invent.\ Math.}
  120; Gritsenko 1999 universal weight) independent of chain fusion. The
  reconstruction theorem above asserts the stronger fusion-consequence
  statement; the level-4-internal identity is unconditional.
\end{itemize}

\section{The $d \geq 4$ stratum}
\label{sec:open-d-geq-4}

\paragraph{Coordinates.} Level 0--4 at $d \in \{4, 5\}$. Chart:
Pandharipande--Thomas at $d = 4$, Fake Monster at $d = 5$. Ambient:
PTVV shifts $-2,-3$ on moduli functions, with Hochschild levels
$E_4,E_5$ before curve specialisation.

\paragraph{Hypotheses.}
\begin{itemize}[nosep]
\item Shift law $(d,\mathrm{PTVV\ shift},\mathrm{function\ bracket},\mathrm{Hochschild})$:
  $(4,-2,\deg 2,E_4)$, $(5,-3,\deg 3,E_5)$.
\item Pandharipande--Thomas DT$_4$ theory at $d = 4$
  (Cao--Maulik--Toda) with virtual cycles in $H_*(\cM)$.
\item Dunn--Lurie additivity $E_5 \simeq E_2 \otimes E_2 \otimes E_1$
  for the Fake-Monster locus $K3 \times K3 \times E$.
\item Borcherds 1990 Fake Monster Lie algebra on $\mathrm{II}_{25, 1}$
  with $\Phi_{12}$ as the universal denominator.
\end{itemize}

\paragraph{Conjecture.} The two-stage
factorisation $\Phi_d = \SpCh_{\Sigma_{d-1}, C} \circ \PhiFA_d$ extends
to $d = 4$ and $d = 5$ with the codomain $E_{n(d)}$ at level 2 forced
by factorisation homology: $E_1$ after specialisation for both
$d=4$ and $d=5$. Before specialisation the Hochschild levels are
$E_4$ and $E_5$, while the PTVV shifts on moduli functions are
$-2$ and $-3$. The chain fusion (\S\ref{sec:open-chain-fusion})
extends correspondingly. The Fake Monster $\mathfrak{g}_{\Phi_{12}}$
on $\mathrm{II}_{25, 1}$ is the desired comparison object for
$\Phi^{(\Sigma_4, C)}_5(D^b\mathrm{Coh}(K3 \times K3 \times E))$.

\paragraph{Constructed loci.} Pandharipande--Thomas DT$_4$ at $d = 4$ for
local $\mathbb{P}^2 \times \mathbb{P}^1$ via Cao--Maulik--Toda; the
shift-law row is consistent. The $K3 \times K3 \times E$ Hodge diamond
$h^{0, q}(X) = (1, 1, 2, 2, 1, 1)$ admits the K\"unneth Stage-1 datum
(Construction~\ref{constr:phi-5-family}); the Stage-2 specialisation
landing on $\mathfrak{g}_{\Phi_{12}}$ is open.

\paragraph{Comparison data.}
\begin{enumerate}[nosep]
\item The Stage-$1$ operad factorises as
  $E_5 \simeq E_2 \otimes E_2 \otimes E_1$ on
  $K3 \times K3 \times E$ by Dunn--Lurie additivity; the
  $\SpCh_{\Sigma_4,C}$ integration then leaves an $E_1$ curve algebra.
\item Construct the boundary vacuum $b(K3 \times K3 \times E, \Sigma_4, C)$
  on the lattice-polarised period domain
  $\Omega^+(\mathrm{II}_{26, 2}) / O^+(\mathrm{II}_{25, 1})$.
\item Verify the universal Borcherds-weight identity at $\Phi_{12}$:
  $\kappa_{\mathrm{BKM}}(\Phi_{12}) = 24/2 = 12$.
\item For $d = 4$, compare the $E_4$ Hochschild object and its
  $E_1$ curve shadow with Pandharipande--Thomas virtual-cycle theory
  and the Conway/Leech bridge.
\end{enumerate}

\paragraph{Obstruction.}
\begin{itemize}[nosep]
\item \emph{Rank obstruction.} The K3-transverse CY$_3$ lane supplies
  the algebraic rank bound \(h^{1,1}(K3)+h^{1,1}(E)=21<24\), so it
  cannot supply the Leech lattice required by
  \(\mathrm{II}_{25,1}\). This is a chart-level obstruction, not a
  symmetric-signature computation on the full \(K3\times E\) threefold
  and not a classification of all CY$_3$ categories.
\item \emph{Equivalence check.} The Conway / Leech datum at $d = 4$ is
  the bridge between the rank-$3$ K3 BKM and the rank-$24$ Fake
  Monster; the Conway lattice $\Lambda_{24}$ embeds in
  $\mathrm{II}_{25, 1}$ along the Leech direction.
\end{itemize}

\section{Higher-$n$ bar = twisting at $E_n$}
\label{sec:open-higher-n-bar}

\paragraph{Coordinates.} Level 2 (bar/twisting) at $E_n$ for $n \geq 2$.
Chart: $E_n$-cooperad / cobar duality. Ambient:
$(\infty, 1)$-categorical bar/cobar adjunction.

\paragraph{Hypotheses.}
\begin{itemize}[nosep]
\item Bar/cobar adjunction at $n = 1$: classical
  (Quillen, Lefèvre-Hasegawa, Loday--Vallette).
\item Higher-$n$ bar construction $B_n(A)$ for $E_n$-algebras
  (Francis--Gaitsgory factorisation co-algebra).
\item $E_n$-cooperad $E_n^{\vee}$ for $n \geq 2$ in the
  $(\infty, 1)$-categorical sense.
\end{itemize}

\paragraph{Conjecture.} For $n \geq 2$, the
bar $B_n(A)$ of an $E_n$-algebra computes twisting/coupling data and
$\Omega_n B_n(A) \simeq A$ in the appropriate ambient. The
identification of $B_n$ with the twisting/coupling data extends the
$n = 1$ classical statement to higher arity.

\paragraph{Constructed loci.} $n = 1$ classical. $n = 2$ Francis--Gaitsgory
factorisation. $n \geq 3$ open.

\paragraph{Construction.}
\begin{enumerate}[nosep]
\item Use the Francis--Gaitsgory factorisation cooperad to construct
  $B_n$ at $n = 2$, identifying it with the $E_2$-twisting data on
  $\Rep^{E_1}(A)$.
\item For $n \geq 3$, use Lurie's higher centralizer
  $\mathfrak{Z}_{E_n}(A)$ to identify the higher bar with the
  higher-twisting data: $B_n(A) = \mathfrak{Z}_{E_n}(A)$ via
  Lurie HA 5.3.1.
\item Verify $\Omega_n B_n(A) \simeq A$ at the cooperad level via the
  Koszul self-duality of $E_n$ for $n \geq 1$ (Fresse).
\end{enumerate}

\paragraph{Obstruction.}
\begin{itemize}[nosep]
\item \emph{Equivalence check.} ``$\simeq$'' for $n \geq 2$ requires
  the appropriate $E_n$-monoidal $\infty$-category; the chain-level
  identification at higher $n$ is involved.
\item \emph{Ambient check.} The Francis--Gaitsgory cooperad lives in
  the derived $\infty$-categorical ambient; the chain-level
  identification requires the Beilinson--Drinfeld factorisation
  $\Ran(X)$-realisation.
\end{itemize}

\section{CY-B$_3$ Koszul, CY-C non-abelian Yangian, CY-D dimensional
stratification at odd $d$}
\label{sec:open-cybcd}

\paragraph{Coordinates.} Level 3 (centre) at $d = 3$ (CY-B, CY-C) and
level 2 (chart-shadow) at odd $d$ (CY-D).

\subsection*{CY-B$_3$ Koszul}

\paragraph{Hypotheses.} The CY-A$_3$ object-level + $E_1$-rigidity
theorem (the working notes) plus the Koszul
self-duality $A^! \simeq A$ on the CY$_3$ datum.

\paragraph{Conjecture.} CY$_3$ smooth proper
admits Koszul self-duality $A^! \simeq A$ at the Stage-1 chiral level.
The $\Omega B_1(A) \simeq A$ comparison holds with explicit denominators
from the CY trace.

\paragraph{Construction.} Use the Beilinson--Ginzburg--Soergel Koszul
duality framework on the chiral side, with the CY trace providing the
explicit perfect pairing. Verify on the Fermat quintic with
$\mathbb{Z}_5^5$-equivariance reducing the Koszul data to a finite
computation.

\subsection*{CY-C non-abelian Yangian}

\paragraph{Hypotheses.} The K3 abelian Yangian (24 Heisenberg
generators with Mukai-signature Serre relations) extended by
non-abelian generators sourced by the Mukai self-mirror branch.

\paragraph{Conjecture.} On $K3 \times E$ the chiral
non-abelian extension of the K3 Yangian
$Y(\mathfrak{g}_{K3})^{\mathrm{non-ab}}$ admits a matrix Miura presentation
and satisfies the $\mathfrak{sl}_2$-Serre constraint $P_2 = 0$ exactly,
with $P_2(D) = \sum_\alpha \alpha^2(k + c_\alpha / D) = 0$ at $D = 3$.

\paragraph{Constructed loci.} Abelian K3 Yangian (24 Heisenberg
generators) chain-level. Matrix Miura on
$V_k(\widehat{\mathfrak{sl}}_2)$. Leading-order $P_2 = 0$ proved;
exact-vanishing conjectural.

\paragraph{Construction.} Verify $P_2 = 0$ exact via Molev's reflection
algebra computation on the $(4, 20)$-signature Mukai datum, using the
Arnaudon--Cramp\'e--Doikou--Frappat--Ragoucy reflection equation. The
target Yangian is $Y_\hbar(\mathfrak{so}(4, 20))$ in the ungraded form
(not $Y_\hbar(\mathfrak{osp}(4 \mid 20))$; the Mukai pairing is symmetric
on both parts, never symplectic on the odd part). The Hodge-parity
super-extension $Y_\hbar(\mathfrak{so}(4 \mid 20))$ (non-Kac) is
available when the Hodge $\mathbb{Z}/2$-super is imposed.

\subsection*{CY-D dimensional stratification at odd $d$}

\paragraph{Theorem.} For compact CY$_d$ with $d$ odd,
\[
 \kappa_{\mathrm{ch}}^{\mathrm{Hodge}}(A_X)
 =
 \sum_q(-1)^q h^{0,q}(X)
 =
0 .
\]
Indeed Serre duality gives $h^{0,q}(X)=h^{0,d-q}(X)$, while the signs
of the paired summands are opposite for odd $d$. Thus the compact Hodge
supertrace cannot be a non-zero Borcherds weight.

\paragraph{BKM scalar.} The BKM value remains a level-$4$ automorphic
scalar:
\[
 \kappa_{\mathrm{BKM}}(\Phi_N)=c_N(0)/2
\]
for the chosen Borcherds input. At $d=3$ and $X=K3\times E$, this gives
$\kappa_{\mathrm{BKM}}(\Delta_5)=5$, while
$\kappa_{\mathrm{ch}}^{\mathrm{Hodge}}(K3\times E)=0$. At $d=5$ the
Fake-Monster input gives $\kappa_{\mathrm{BKM}}(\Phi_{12})=12$, again
independent of the compact Hodge supertrace.

\section{Non-abelian K3 Yangian for non-simply-laced $\fg_{K3}$}
\label{sec:open-nonabelian-K3-Yangian}

\paragraph{Coordinates.} Level 3 at $d = 2$ on the K3 Mukai self-mirror
branch with non-simply-laced $\fg_{K3}$. Chart: K3 with
non-simply-laced ADE-like structure on the Mukai lattice. Ambient:
chain-level Yangian.

\paragraph{Hypotheses.}
\begin{itemize}[nosep]
\item K3 abelian Yangian on the Mukai datum
  (24 Heisenberg generators).
\item Costello--Gaiotto--Yagi 5d hCS construction at simply-laced
  $\fg$: yields the Yangian VOA $Y^{\mathrm{VOA}}(\fg)$ to all orders
  in $\hbar$.
\item Non-simply-laced extension: the twisted Yangian
  $Y^{\mathrm{tw}}(\fg)$ on a non-simply-laced ADE-like
  structure on the Mukai lattice.
\end{itemize}

\paragraph{Conjecture.} Non-simply-laced
$\fg_{K3}$ admits a twisted Yangian
$Y^{\mathrm{tw}}_\hbar(\fg_{K3})$ on the K3 chiral algebra at level 3,
quantising the 5d hCS construction in the twisted sector.

\paragraph{Constructed loci.} Simply-laced $\fg$ via Costello--Gaiotto--Yagi.
Non-simply-laced is open at all orders in $\hbar$.

\paragraph{Construction.}
\begin{enumerate}[nosep]
\item Identify the non-simply-laced ADE-like structure on the Mukai
  lattice $H^{\mathrm{ev}}(K3, \mathbb{Z})$ via the
  Eichler--Zagier correspondence.
\item Apply the Molev twisted-Yangian construction to the Mukai-pairing
  data, with the Chevalley involution $\theta$ implementing the twist.
\item Extend the 5d hCS computation to the twisted sector via the
  $\theta$-equivariant BV propagator on
  $\mathbb{R} \times \mathbb{C}^2$.
\end{enumerate}

\paragraph{Obstruction.}
\begin{itemize}[nosep]
\item \emph{Theory-import check.} The 5d hCS construction is
  known at simply-laced $\fg$ by Costello--Yagi, not by the false
  assertion that $d^{abc}$ vanishes on every ADE algebra. Cubic
  anomaly freedom is the separate invariant-theoretic condition
  excluding $A_n$, $n \geq 2$; the non-simply-laced 5d problem is open
  because the all-orders twisted comparison is not yet constructed.
\end{itemize}

\section{CY-A$_3$ chain-level compact closure}
\label{sec:open-cya3-chain}

\paragraph{Coordinates.} Level 1 (Stage-1 chain-level) at $d = 3$ for
compact non-formal CY$_3$. Chart: compact non-toric CY$_3$ (quintic,
generic Schoen). Ambient: chain-level with $\Ainf$-compatible
$S^3$-framing on $HC^-_3(C)$.

\paragraph{Hypotheses.}
\begin{itemize}[nosep]
\item CY-A$_3$ object-level + $E_1$-rigidity
  ($\infty$-categorical level proved).
\item Chain-level TCFT datum: $E_3$-formality point,
  $\Ainf$-compatible $S^3$-framing homotopy on $HC^-_3(C)$, anomaly
  cancellation, analytic completion, admissible Stage-2
  specialisation.
\end{itemize}

\paragraph{Conjecture.} The chain-level $\Ainf$
two-stage factorisation
$\Phi_3 = \SpCh_{\Sigma_2, C} \circ \PhiFA_3$ extends to compact
non-formal CY$_3$ in the J-adic and pro-object ambients of the raw
bar complex.

\paragraph{Constructed loci.} Constructed loci (toric $\mathbb{C}^3$, local
$\mathbb{P}^2$, resolved conifold, $K3 \times E$). Compact non-formal
quintic and generic Schoen are open.

\paragraph{Construction.}
\begin{enumerate}[nosep]
\item Construct the $\Ainf$-compatible $S^3$-framing on the Fermat
  quintic via the $\mathbb{Z}_5^5$-equivariant cover, reducing the
  computation to invariant-sector $\Ainf$-data.
\item Verify anomaly cancellation via Costello--Gwilliam--Li one-loop
  BV: the cubic-symmetric-Casimir $d^{abc}$ must vanish for the
  chosen Stage-1 gauge input; simple-lacing does not imply this.
\item Extend the framing through the J-adic completion to the
  non-formal locus, with the direct-sum chain-level model included in
  the datum.
\end{enumerate}

\paragraph{Obstruction.}
\begin{itemize}[nosep]
\item \emph{False-functoriality check.} Strict formality of
  $A_{\mathrm{BVDB}} = \mathrm{End}^\bullet(\bigoplus_{i=0}^4 \cO_{X_5}(i))$
  on the quintic is refuted: Yukawa coupling $Y_3 = H^3 = 5$ is the
  non-zero Kodaira--Spencer $m_3$. Correct statement: curved
  formality with Yukawa as BCOV datum.
\item \emph{Numerical-constant check.} The number $204$ is the Hodge
  diamond total $1 + 101 + 101 + 1$, not a $\mathbb{Z}_5^5$-invariant
  sector dimension; computations must use the correct invariant-sector
  count.
\end{itemize}

\section{Stage-1 four-physical-lane completeness}
\label{sec:open-four-lanes}

\paragraph{Coordinates.} Level 1 (Stage-1 physical realisations) at
$d \in \{2, 3\}$. Chart: physical-realisation lane (5d hCS, 6d hCS,
mixed-HT-strings local model, Costello--Gwilliam--Li perturbative).
Ambient: BV cohomology + holomorphic locality.

\paragraph{Hypotheses.}
\begin{itemize}[nosep]
\item 5d hCS on $\mathbb{R} \times \mathbb{C}^2$:
  Costello--Gaiotto--Yagi, all orders in $\hbar$ for simply-laced $\fg$.
\item 6d hCS on $\mathbb{C}^3$: Costello--Li, classical / finite-Rees
  level constructed; one-loop anomaly $\int_X \mathrm{Tr}_{\mathrm{ad}}(A(F_A)^3)$
  sourced by cubic Casimir $d^{abc}$, vanishing on $\fsl_2$,
  obstructing $\fsl_{N \geq 3}$ with $d^{abc}d_{abc} = 2N$ in the chosen
  trace normalisation.
\item Mixed-HT-strings local model: $\mathbb{R}^2_{\mathrm{top}} \times \mathbb{C}^2_{\mathrm{hol}}$
  with Hamiltonian BF sector; the holomorphic de Rham obstruction
  controls the global lift on non-formal targets.
\item Costello--Gwilliam--Li perturbative: factorisation algebras of
  classical / quantum observables.
\end{itemize}

\paragraph{Conjecture.} The four physical lanes of
Stage-1 produce equivalent $E_d$-holomorphic factorisation algebras at
the constructed loci, with the chain-level and \(\infty,1\)-categorical constructions carrying equal status. The compact non-toric extension lives along the holomorphic de
Rham class on the mixed-HT-strings lane.

\paragraph{Constructed loci.} 5d hCS $\to$ Yangian VOA at $d = 2$
simply-laced. 6d hCS on $\mathbb{C}^3$ at $\fsl_2$ unobstructed; resolved
conifold two-chart descent constructed at finite-Rees level with
anomaly-free perturbative quantum lift. Mixed-HT-strings local model
verified at toric loci. Costello--Gwilliam--Li perturbative on
$\mathbb{C}^d$ all-orders.

\paragraph{Construction.}
\begin{enumerate}[nosep]
\item Twisted Yangian construction for non-simply-laced $\fg$ in 5d
  hCS at all orders (\S\ref{sec:open-nonabelian-K3-Yangian}).
\item Compact $K3 \times E$ Hall double via the 6d hCS lane: conditional on
  reduced compact Hall data, negative-half / pairing / radical
  recognition, and CY-A$_3$ chain-level closure
  (\S\ref{sec:open-cya3-chain}).
\item Mixed-HT-strings holomorphic de Rham obstruction on compact
  non-toric: identify the obstruction class as the level-1 $\to$ level-1
  obstruction to the global lift, computed via the
  Costello--Li one-loop BV residue.
\end{enumerate}

\paragraph{Obstruction.}
\begin{itemize}[nosep]
\item \emph{Theory-import check.} 6d hCS is not 3d Chern--Simons; the
  one-loop obstruction is quartic in fields, not cubic. The cubic Casimir
  $d^{abc}$ contributes the cohomological obstruction. The BV-trivial
  counter-term absorbs only the wave-function piece
  $A_{\mathrm{w.f.}} = -C_2 / (2\pi)^3 = -2 h^\vee / (2\pi)^3$.
\item \emph{Ambient check.} The Miki $S_3$ torus-Weyl statement has no
  K3 torus-action input; only the elliptic-factor $\mathrm{SL}_2(\mathbb{Z})$
  survives on $K3 \times E$.
\end{itemize}

\section{Scope-determination obligations carried by other parts}
\label{sec:open-scope-obligations}

The discipline applied to the body of the volume gives a finite list of
coordinate conditions for theorem statements in which at least one axis is
underdeclared. Each condition is mathematical: without the missing
coordinate the statement is not well typed.

\begin{itemize}[nosep]
\item \emph{Stage-2 chart-internal disambiguation.}
  Chapter~\ref{ch:k3e-bkm} constructs six $(\Sigma_2, C)$-routes
  to $G(K3 \times E)$. Each route's chart-internal coordinates
  (boundary vacuum $b$, admissibility window) are named per route, but
  the explicit chart label per row of the comparison table requires
	  verification in the comparison table.
\item \emph{Class~$\cM$ bar-tower ambient.} The class~$\cM$ $E_3$
  bar identity $=6^g$ is a weight-completed statement; ordinary chain
  complexes do not carry the finite tower in this form.
\item \emph{Six routes versus six $\Phi$-applications.} The structural
  remark that the six K3 $\times E$ routes are six
  $(\Sigma_2, C)$-specialisations of one Stage-1 datum
  $\PhiFA_3(D^b\mathrm{Coh}(K3 \times E))$, not six $\Phi$-applications,
	  belongs in every row of the K3 $\times E$ catalogue.
\item \emph{$\kappa_{\mathrm{ch}}$ subscript discipline.} Every
  body occurrence of $\kappa_{\mathrm{ch}}$ carries a reading
  superscript identifying its construction layer. The subscripted
  variants in active use:
  $\kappa_{\mathrm{ch}}^{\mathrm{Hodge}}, \kappa_{\mathrm{ch}}^{\mathrm{Heis}},
  \kappa_{\mathrm{ch}}^{\mathrm{Mukai}}, \kappa_{\mathrm{ch}}^{\mathrm{cpt}},
  \kappa_{\mathrm{ch}}^{\mathrm{loc}}, \kappa_{\mathrm{ch}, \mathrm{BV}}$.
  Every body occurrence carries the relevant subscript.
\item \emph{Bare $\Phi$ specialisation.} The naked symbol $\Phi$ at
  body level decomposes as either
  $\Phi^{(\Sigma_{d-1}, C)}_d$ (Stage-2 specialised) or $\PhiFA_d$
  (Stage-1 native). Every occurrence carries the relevant superscript.
\item \emph{Bulk versus bar.} The universal closed sector is
  $\Zder(A)$ (the derived chiral centre at level $3$); a physical
  bulk identification requires the OCA arrow
  $\cO^{\mathrm{phys}}_{\mathrm{bulk}}(\mathcal T) \xrightarrow{\sim} \Zder(A)$
  of Definition~\ref{def:oca-quasi-iso}. The bar $B(A)$ at level $2$
  is twisting/coupling, not closed-sector; the comparison
  $\Zder(A) = \RHom(\Omega B(A), A)$ is the named arrow.
\end{itemize}

The list does not add new mathematics; it records the coordinates required
for the stated comparison maps to be well typed.
